\section{Introduction}
Mathematical knowledge inherently forms a vast network of logical dependencies, an architecture that has historically remained unrecorded. Following long-standing convention, traditional mathematical writing preserves final results while heavily compressing the underlying logic, frequently omitting the tacit deductions and canonical foundations that support the work. Because the density of these explicit details is dictated by publication constraints and authorial stylistic preferences, the global web of logical relations exists only as a latent, epistemically opaque construct, obscured by coarse-grained citations. The advent of formal proof assistants fundamentally alters this paradigm, converting an implicit, community-held mental model into an explicit, machine-analyzable topology.

Beginning with de~Bruijn's Automath~\cite{debruijn1970}, and followed by Mizar~\cite{trybulec1978}, Rocq/Coq~\cite{coquand1988}, and Isabelle~\cite{paulson1994}, these systems employ a verification kernel to check every declaration and record the precise premises it invokes~\cite{avigad2024}. Building on this tradition, \module{Mathlib}, the mathematical library for Lean~4~\cite{mathlib2020,lean4}, is a natural candidate to express logical dependencies and it has grown to a scale where macroscopic statistical structure might emerge.
Ultimately, the value of \module{Mathlib} lies not merely in aggregating verified truths, but in the structural discoverability it affords. As Feng et al.~\cite{feng2026erdos} observe in their study of the Erd\H{o}s problems, today the difficulty of solving open conjectures often lies not in the mathematics itself, but in the obscurity of the existing literature. Libraries like \module{Mathlib} have the potential to mitigate this issue by transforming scattered results into a unified and queryable structure. That said, it is important to observe that while the validity of each theorem is enforced deterministically by the kernel, the library's macroscopic organization emerges from a human-driven process. This includes naming conventions inherited from centuries of mathematical practice, file hierarchies shaped by collaborative workflows, and namespace boundaries drawn by cognitive habit. Consequently, these network properties do more than document logical dependencies; they reveal measurable traces of how mathematics is produced, reflecting community behavior and the nuances of language design.

The present paper analyzes these structural patterns in the attempt to quantify the divergence between inherited mathematical structure (the product) and human organizational process (the process). Our goal is to establish a baseline for human-authored formal mathematics, particularly as AI-driven formalization and automated discovery increasingly capture the attention of the mathematical community.~\cite{hubert2025alphaproof,achim2025aristotle,chen2025seedprover,xin2025deepseekprover}. We take \module{Mathlib} as the central object of study: 308,129 declarations, 8,436,366 dependency edges, and 7,563 modules. All statistics refer to commit \texttt{534cf0b} (2 Feb 2026), extracted using \texttt{importGraph} (module layer) and premise-extraction tooling (declaration layer). Our module import graph reflects a post-\texttt{shake} state: Mathlib's CI runs the \texttt{shake} linter, which automatically removes certain transitively redundant imports; the raw developer-written imports would contain more redundancy. Moreover, the Lean~4 module system (introduced November 2025) distinguishes \texttt{public import} from \texttt{import}~\cite{lean4modules}; our snapshot captures a transitional state where most imports are still \texttt{public}, and the import graph will change structurally as the community gradually adopts private imports. We develop our analysis at four progressively finer levels of granularity:

\medskip
\noindent\textbf{The module graph (\S\ref{sec:module-import}).}
Each module (source file) is a node. The \texttt{import} statements form edges. We analyze two graphs simultaneously together with the systematic divergence between them: the module tree~$T$, derived from the dot-separated hierarchy of modules (\S\ref{sec:namespace-tree}), which encodes the human cognitive classification of mathematics; the module graph~$G_{\mathrm{module}}$, which encodes the compiler-verified logical dependencies. This granularity captures, at least in part, the \emph{architecture} of the library, \emph{i.e.}, its overall shape, organizational principles, and structural vulnerabilities.

\medskip
\noindent\textbf{The declaration graph (\S\ref{sec:theorem-premise}).}
Each theorem or definition is a node. The premises invoked in a proof form edges. This level penetrates into the mathematical content itself---no longer dependencies between files, but logical relations between theorems. It measures the \emph{reasoning structure} of mathematics, \emph{i.e.}, which theorems serve as core lemmas, which are isolated special cases, and how knowledge flows across theories.

\medskip
\noindent\textbf{The namespace graph (\S\ref{sec:namespace-graph}).}
It offers an intermediate granularity derived from the declaration graph. Each declaration's fully qualified name defines a hierarchy of namespaces at varying depths; truncating to depth~$k$ and aggregating declaration-level edges yields a family of dependency graphs $G_{\mathrm{ns}}^{(k)}$. We measure how \emph{containment} (Definition~\ref{def:containment})---the proportion of edges staying within namespace boundaries---decays from $48.9\%$ at the top-level directory to ${\sim}13\%$ at finer namespaces. This decay curve quantifies the scale at which human cognitive classification retains organizing force and the scale at which mathematical logic overwhelms it. At the coarsest level, broad mathematical disciplines capture nearly half of all dependencies; at finer levels, most dependencies escape any single namespace.

\medskip
\noindent\textbf{Cross-level analysis (\S\ref{sec:module-declaration}).}
The preceding three levels treat modules, namespaces, and declarations as separate objects. In \S\ref{sec:module-declaration} we overlay them so that each module becomes a \emph{region} in the declaration dependency graph, and we measure how well the human-drawn module and namespace boundaries align with the logical dependency structure. Module cohesion (Definition~\ref{def:cohesion}), or, more explicitly, the proportion of a module's dependency traffic that stays internal, measures how far reality departs from the ideal in which module boundaries capture logical structure.

\medskip
Together, the four levels form a progression from coarse architecture (\S\ref{sec:module-import}) through fine-grained logical structure (\S\ref{sec:theorem-premise}) and its namespace-level aggregation (\S\ref{sec:namespace-graph}) to the tension between human organization and 'mathematical necessity' (\S\ref{sec:module-declaration}). Beyond characterizing the network, we identify structural metrics---module cohesion, centrality rankings, zero-citation rates, containment curves---that can serve as quantitative signals for library maintenance, compilation optimization, and AI training data curation.

\subsection{Previous works}
\paragraph{Software dependency networks.}
Dependency graphs arise naturally in software ecosystems: package managers such as npm, PyPI, and Maven define directed acyclic graphs in which each package declares its requirements~\cite{decan2019,labelle2004}. In the formal mathematics setting, Blanchette et al.~\cite{blanchette2015} mine structural properties of the Isabelle AFP, making it the closest methodological predecessor to our work, and Bancerek et al.~\cite{bancerek2018} document the MML's role and evolution. Studies of these ecosystems have revealed heavy-tailed degree distributions, small-world connectivity, and fragility under targeted removal of hub packages. Our analysis applies a similar methodology to a formal mathematics library---a setting in which every edge carries a stronger guarantee than ``package~$A$ requires package~$B$ at build time.'' In a type-checked library, premise-level dependencies are mechanically certified: if a constant appears in a proof term, the kernel enforces that dependency. Module-level imports, by contrast, are developer-declared and can be redundant relative to reachability ($17.5\%$ are transitively implied, measured after Mathlib's \texttt{shake} linter has already pruned certain redundant imports in CI). In our opinion, the graph records not merely a build requirement but a semantic relationship.

\paragraph{\module{Mathlib} growth and maintenance.}
Van Doorn, Ebner, and Lewis~\cite{vandoorn2020} described the engineering practices that sustain \module{Mathlib} as a collaboratively maintained library, addressing linting, documentation, and refactoring workflows. Ringer et al.~\cite{ringer2019} provide a canonical survey of proof engineering, noting that ``there is still relatively little work describing tools and best practices''; our contribution fills part of this gap. Commelin and Topaz~\cite{commelin2023} proposed abstraction boundaries and specification-driven development as organizing principles for formal mathematics. More recently, Commelin et al.~\cite{baanen2025} documented the challenges of scaling \module{Mathlib}, including build times, namespace reorganization, and technical debt. On the tooling side, Massot's \textsc{leanblueprint} system~\cite{massot2020leanblueprint} provides infrastructure for managing large formalization projects: mathematicians manually annotate dependency relations between statements in \LaTeX{} using \verb|\uses{}| tags, producing a dependency graph that tracks formalization progress. This approach has been adopted by several major projects, including the Liquid Tensor Experiment and the ongoing formalization of Fermat's Last Theorem. Zhu et al.~\cite{zhu2025leanarchitect} complement this with \textsc{LeanArchitect}, which automatically infers dependency information from Lean source code, eliminating the manual duplication between informal and formal representations. Notably, their work includes AI integration: in a case study, the automatically extracted dependency graph guides an AI prover (Achim et al.'s Aristotle system~\cite{achim2025aristotle}) to formalize theorems in topological order. Our work differs from both in scope and purpose: \textsc{leanblueprint} and \textsc{LeanArchitect} extract dependency graphs to manage individual formalization projects; we analyze the global dependency structure of the entire \module{Mathlib} library as a network, applying tools from network science---degree distributions, community detection, robustness analysis---to characterize structural patterns that no single project would reveal. Avigad et al.~\cite{avigad2025anatomy} describe the practitioner experience of using proof assistants; Klein et al.~\cite{klein2012} address managing large-scale proofs in Isabelle, including tools for moving lemmas across modules. These works address \module{Mathlib}'s maintenance from a software engineering perspective. Our analysis complements theirs by providing a network-theoretic diagnostic of the same challenges they address qualitatively: we show, for example, that the $17.5\%$ import redundancy rate (Definition~\ref{def:redundancy-rate}, measured post-\texttt{shake}) and $0.107$ module cohesion quantify the development ergonomics overhead that maintenance practices must manage.

\paragraph{Machine learning for formal mathematics.}
The graph structure of \module{Mathlib} has been exploited in machine learning, particularly for premise selection. Yang et al.~\cite{leandojo} introduced the LeanDojo benchmark, extracting proof traces from \module{Mathlib} that implicitly define a dependency graph. Bemberek et al.~\cite{bemberek2025} constructed a heterogeneous dependency graph from this benchmark and applied relational graph convolutional networks for premise selection, achieving a $25\%$ improvement over text-only baselines. Lu et al.~\cite{leanfinder} built a directed acyclic graph of Lean~4 statements with dependency edges to support semantic search. On the knowledge-representation side, the EpiK Protocol~\cite{epik2024} constructed a knowledge graph from Lean~4 and \module{Mathlib}, applying embedding models for link prediction. Uskuplu and Moss~\cite{knowtex2025} argued that dependency graphs should become a standard feature of mathematical writing. These systems consume the dependency graph as training data or as a retrieval index. We study the graph itself---not as input to a downstream task, but as an object of independent structural interest.

\paragraph{Complex network methodology.}
The analytical tools we employ are basic and standard in network science. Heavy-tailed network theory~\cite{barabasi1999}, small-world analysis~\cite{watts1998}, and modularity-based community detection~\cite{newman2004,blondel2008} have been applied to biological, social, and technological networks, but not, to our knowledge, to a large formal mathematics library at the multi-layer depth we pursue here. (Blanchette et al.~\cite{blanchette2015} apply mining techniques to the Isabelle AFP, and Huch and Wenzel~\cite{huch2024} analyze the AFP's theory DAG for parallel build scheduling, but neither pursues the three-layer network analysis we develop.) In a related setting, Barbosa et al.~\cite{barbosa2013} studied the network structure of three online mathematical libraries---Wikipedia (restricted to mathematics), MathWorld, and DLMF---finding large strongly connected components, the absence of clear power laws, and the effectiveness of stress centrality for navigating mathematical content. Their complex networks, however, are hyperlinked encyclopedias with informal content, while here we consider a machine-verified library in which every dependency is logically certified.

\subsection{Our Contribution}
\label{sec:contribution}
We mathematically define three graph objects from \module{Mathlib} and analyze them as a multi-layer dependency network:
\begin{itemize}
    \item the module import graph~$G_{\mathrm{module}}$ (Definition~\ref{def:module-graph}), a DAG whose nodes are source files and whose edges are \texttt{import} statements;
    \item the declaration dependency graph~$G_{\mathrm{thm}}$ (Definition~\ref{def:thm-graph}), a DAG whose nodes are theorems and definitions and whose edges are premise invocations;
    \item the namespace graphs~$G_{\mathrm{ns}}^{(k)}$ (Definition~\ref{def:ns-graph}), obtained by aggregating declaration-level edges at varying namespace depths;
\end{itemize}
together with the module tree~$T$ derived from the dot-separated file hierarchy (Definition~\ref{def:module-tree}). This formalization treats \module{Mathlib} as a graph-theoretic object, separating module imports and namespace classification from declaration-level dependencies. The former two reflect tooling and collaboration constraints; the latter records the complete logical network of traditional mathematics. This separation enables controlled comparison between human organizational process and inherited mathematical structure.

We provide a quantitative characterization of \module{Mathlib}'s dependency structure across all three layers---degree distributions (\S\ref{sec:sa-degree}), DAG depth profiles (\S\ref{sec:dag-depth}), transitive redundancy (\S\ref{sec:transitive-reduction}), containment curves (\S\ref{sec:containment-decay}), centrality rankings (\S\ref{sec:sa-centrality}), community structure (\S\ref{sec:sa-community}), and robustness under node removal (\S\ref{sec:sa-robustness})---aiming to establish a preliminary empirical baseline (\S\ref{sec:structural-analysis}) and derive actionable implications for library maintenance, AI-driven formalization, and tooling (\S\ref{sec:conclusion}).

Comparing layers reveals that \textbf{organization diverges from dependency}: the module and namespace communities align with each other (NMI\,$=$\,0.71; Definition~\ref{def:nmi}, \S\ref{sec:sa-community}) but poorly with the declaration-level structure (NMI\,$=$\,0.34). The module layer compresses ${\sim}40\times$ in nodes and ${\sim}360\times$ in edges relative to the declaration layer (\S\ref{sec:module-declaration}), acting as a compressed interface shaped by tooling, while declaration-level dependencies encode the organizational habits of traditional mathematics under the constraint of logical consistency. At the same time, \textbf{17.5\% of module imports are logically redundant} yet structurally essential (\S\ref{sec:transitive-reduction}), even after Mathlib's \texttt{shake} linter has pruned certain redundant imports; these links serve development ergonomics by easing refactoring and ensuring transitive visibility for future extensions.

The three layers also \textbf{exhibit markedly different topologies}: the module graph forms a 153-layer deep funnel (\S\ref{sec:dag-depth}), the declaration graph is shallow and top-heavy (84 layers, 38.8\% at layer~0; \S\ref{sec:thm-dag-depth}), and the namespace graph is nearly flat (8 layers after SCC condensation; \S\ref{sec:ns-dag-depth}). This inversion---human-organized structure being deepest, logical structure being shallowest---is compounded by a flattening of traditional semantic distinctions: formalization preserves labels like `theorem' and `lemma' but erodes their meaning, while automation tools inflate the library with mirrored results sharing identical logical structure (\S\ref{sec:thm-vs-lemma}).

Based on these findings, we discuss concrete implications for library maintenance, compilation optimization, language design, and AI training data curation in \S\ref{sec:practice}.
